% Blue draft inserts for the annual full price-path RE revision.
% These are LaTeX-only drafting notes. Do not paste into the LyX source until approved.
% Assumption for this draft: the full price-path RE runs pass at the required horizons.

% ---------------------------------------------------------------------------
% Insert after the paragraph that currently states the random-walk expectation assumption.
% ---------------------------------------------------------------------------

{\color{blue}
In the revised transition exercises, we separate the original current-price expectation from a full price-path rational expectations exercise. The current-price version keeps the assumption that households do not forecast future demographic changes in house prices. It remains useful as a direct comparison with the original model and as a diagnostic for the political supply channel. The rational expectations version instead treats the future demographic path as known and solves households' choices on the anticipated future path of house prices. The equilibrium price path is therefore a fixed point: a guessed future price path determines household choices, ownership, and political pressure, and this political pressure generates the price path that households must have anticipated.

This extension is computationally more demanding because the demographic state is high dimensional. The economy contains annual age cells from age 25 to age 80, and a demographic shock changes the relative mass of voters at each age as the shocked cohort moves through the life cycle. We therefore use the same annual timing as the calibrated steady state and solve the transition over an 80 year horizon, with additional terminal years when needed so that temporary demographic shocks can dissipate before the terminal condition binds.
}

% ---------------------------------------------------------------------------
% Insert after the smoothed political-pressure definition.
% ---------------------------------------------------------------------------

{\color{blue}
The smoothed political rule is used in both transition exercises. Let $m_t$ denote the aggregate vote imbalance in year $t$, measured relative to the benchmark political target. We summarize the political response by the bounded pressure index
\[
    \pi_t = \tanh(m_t/\tau),
\]
where $\tau$ governs how quickly political pressure rises as the vote imbalance moves away from zero. As $\tau$ becomes small, this rule approaches the hard sign rule in the static median-voter model. For positive $\tau$, political pressure changes continuously with the strength of the vote imbalance. The transition then maps political pressure into the house-price and permit environment using one pass-through parameter, so the dynamic exercise remains tied to the paper's political supply channel without separately identifying several weakly disciplined supply-side parameters.
}

% ---------------------------------------------------------------------------
% Insert in the demographic projection section, after the steady-state projection paragraph.
% ---------------------------------------------------------------------------

{\color{blue}
The steady-state projection remains a useful benchmark because it isolates how the changing age structure shifts the political equilibrium when all other aggregate conditions are held fixed. The transition exercises add a different object. They ask how households behave when they anticipate the future house-price path generated by the demographic projection itself. We therefore report a dynamic counterpart to the steady-state projection under the low, medium, and high immigration paths. In each case, the current-price transition and the full price-path rational expectations transition are plotted on the same axes. The current-price path shows what the original expectation assumption implies, while the rational expectations path shows how much of the future demographic pressure is capitalized into current choices and prices.

In the figures, the steady-state projection is kept as a separate benchmark figure. The dynamic projection figure then uses the same immigration scenarios but compares expectation assumptions within each scenario. Solid lines report the full price-path rational expectations transition; dashed lines report the current-price transition. This keeps the comparison visually disciplined: the reader can see whether rational expectations mainly changes the timing of adjustment, the level of the price response, or both.
}

% ---------------------------------------------------------------------------
% Insert at the start of the baby-boom impulse-response subsection.
% ---------------------------------------------------------------------------

{\color{blue}
For the temporary baby-boom experiment, the key distinction between the two transition exercises is timing. Under the current-price expectation, households respond to the current price and the current political environment, while the realized economy moves forward as the large cohort ages. Under full price-path rational expectations, households also anticipate the future price effects of that cohort as it moves from young renters into middle-aged homeowners and then into older households. This means the future political effect of the boom can be partly capitalized into earlier housing choices and prices.

We therefore report the baby-boom transition in two layers. The main panel plots the current-price and full rational expectations price paths on the same graph. Additional panels report homeownership, housing demand, and political pressure. The hard-vote and random-price exercises are not plotted as additional main lines because they answer a different question. They are diagnostic comparisons and are reported separately, either in the text or in an appendix figure, to show that the benchmark is not driven by arbitrary price movements or by a numerical discontinuity in the voting rule.
}

% ---------------------------------------------------------------------------
% Insert at the start of the permanent demographic shock subsection.
% ---------------------------------------------------------------------------

{\color{blue}
We refer to this second experiment as a permanent decline in entrant formation. This is the dynamic counterpart of the permanent demographic shock in the original paper. The entrant rate is permanently lower, but the effect on the age distribution is transitional because the smaller entering cohorts gradually move through the life cycle. For this reason, the terminal condition differs from the temporary baby-boom experiment. The temporary boom can return to the original demographic environment, while the permanent entrant decline is anchored to the terminal demographic composition implied by the new entrant process.

As with the baby-boom experiment, we plot the current-price transition and the full price-path rational expectations transition on the same axes. If the two paths differ mainly near the start of the transition, this indicates capitalization of future demographic pressure. If they differ mainly near the end, this indicates that the terminal demographic composition is important for the price path. The steady-state results remain useful here because they show the long-run political equilibrium associated with a given demographic structure; the transition results show how the economy moves toward that object when households anticipate the future path.
}

% ---------------------------------------------------------------------------
% Insert in the computational appendix, replacing the current-price-only annual algorithm note.
% ---------------------------------------------------------------------------

{\color{blue}
The annual transition is solved in two versions. The first version keeps the current-price expectation used in the original model. In each year, households solve their problem at the current house price, the distribution is moved forward using the annual transition matrix, and the political pressure generated by the new distribution updates the realized price path.

The second version solves a full price-path rational expectations transition. The algorithm begins with a guessed path of house prices, $\{p^h_t\}_{t=1}^T$. Given this path, the household problem is solved backward, so current choices depend on the current price and continuation values depend on future prices. The distribution is then simulated forward. In each year, the age-weighted political vote generates an implied price path through the smoothed political-pressure rule. A transition equilibrium is a path for which the guessed price path and the implied price path coincide, up to the numerical tolerance used in the computation. For temporary shocks, the internal solution horizon can extend beyond the reported 80 years, so that the shock can unwind before the terminal condition affects the reported path.
}

% ---------------------------------------------------------------------------
% Figure plan for the revised paper. This is not prose unless converted later.
% ---------------------------------------------------------------------------

{\color{blue}
\paragraph{Figure plan for the revised transition section.}
The steady-state demographic figures should remain in the paper because they answer a different question from the transition exercises. They show the political equilibrium associated with each demographic structure. The new transition figures should then compare expectation assumptions, not replace the steady-state figures.

For the baby-boom and permanent entrant-decline experiments, the main transition figure should plot the current-price path and the full price-path rational expectations path on the same axes. The cleanest layout is a multi-panel figure: house price, homeownership, housing demand, and political pressure. The current-price line should be dashed and the full rational expectations line should be solid. For the low, medium, and high immigration projections, use either three panels, one for each immigration assumption, or one main panel for the median projection with low and high projections reported as robustness panels. Random-price paths and the hard-sign $\tau=0$ case should not be added to the main shock graphs unless the figure remains readable. They are better used as diagnostic figures or appendix comparisons.
}
